\section{讨论}
\label{sec:discussion}

\subsection{主要发现总结}
本实验完成了从 sgRNA 构建、慢病毒包装、MC38 细胞感染与分选，到 OT-1 T 细胞共培养杀伤评估的完整流程。结果部分显示：
\begin{enumerate}
    \item LentiV2-puro 载体经 BsmB1 酶切后获得预期线性化片段并成功回收；
    \item Sanger 测序比对提示 sgRNA 插入片段在关键区域与设计序列一致，支持载体构建成功；
    \item 在体外共培养体系中，以 $\log_2\mathrm{FC}$ 表征 KO 相对 NC 的存活变化：总体上，KO-1 与 KO-2 组在 1:6 与 1:9 条件下多呈现 $\log_2\mathrm{FC}>0$ 的趋势，而 NC 对照组整体接近 0，提示 \textit{B2m} 敲除可能降低了 OT-1 T 细胞对靶细胞的相对杀伤效率；但在当前数据规模与方差水平下，组间差异未达到统计显著（图~\ref{fig:log2FC}）。
\end{enumerate}

\subsection{对结果的免疫学解释}
在本实验设计中，CTV$^+$ 群体为对照靶细胞（NC），CTV$^-$ 群体为敲除靶细胞（KO），两者在同孔内共同接受相同的 T 细胞压力与培养条件。若 \textit{B2m} 敲除导致 MHC-I 表达或稳定性下降，则在 OVA 肽脉冲后，敲除靶细胞表面形成并展示 SIINFEKL-H2K\textsuperscript{b} 的能力将受限，从而减少 OT-1 T 细胞的特异性识别与杀伤。该机制预期带来 KO 相对 NC 的保留或富集，在本报告采用的比值框架下表现为 $r$ 相对升高、$\log_2\mathrm{FC}>0$。从多数孔位的方向性来看，这一趋势与 \textit{B2m} 在抗原呈递中的功能一致。但部分样本存在较大幅度的负向 $\log_2\mathrm{FC}$ 或孔间变异（例如个别孔位出现极端比值），整体未获得统计学显著差异。

\subsection{潜在误差来源与局限性}
\begin{enumerate}
  \item \textbf{敲除效率与群体异质性：} 本实验以流式统计的存活比例作为终点，但未同步呈现 \textit{B2m} 蛋白水平（例如 H2K\textsuperscript{b} 表达）或 indel 水平的定量敲除效率。若 KO 群体中存在较高比例的未敲除细胞，将削弱表型差异。
  \item \textbf{混合比例与归一化假设：} 本实验采用同组 only 的均值 $\overline r_{\mathrm{only}}$ 作为基线进行归一化。该方法能够部分抵消初始混样比例偏差，但前提是 only 条件能稳定反映同批次细胞的``无杀伤基线''。若贴壁效率、增殖速率或计数误差导致 only 孔间 $r$ 波动，则会放大 $\log_2\mathrm{FC}$ 的方差。
  \item \textbf{门控与检测噪声：} CTV 信号衰减、补偿设置、碎片排除与双ts细胞（doublet）排除策略均可能影响 CTV$^+$/CTV$^-$ 的门控比例。尤其在 KO/NC 比值接近 1 时，小幅门控漂移即可造成显著的比值波动。
  \item \textbf{E:T 条件与共培养时间窗：} 当前仅测试 1:6 与 1:9 两个 E:T 比率，且以单一时间窗终点读出。若杀伤动力学存在非线性（例如早期快速杀伤后进入平台期），不同条件间差异可能被终点采样低估。
  \item \textbf{统计功效：} 当前每条件为 3 个技术重复孔，且不同同学样本之间存在批次差异。若需进行稳健的组间推断，仍需要增加生物学重复并统一实验批次。
\end{enumerate}

\subsection{改进建议与后续工作}
为提高结论可靠性与可解释性，在后续实验中可以进行如下优化：
\begin{enumerate}
  \item 在共培养前后分别用抗体染色定量 H2K\textsuperscript{b}/B2m 相关表面标志，或采用 TIDE/ICE 等方法对 indel 进行定量，以建立敲除效率—表型强度的对应关系。
  \item 扩展 E:T 梯度（如加入 1:3）并增加时间点（如 8 h、16 h、24 h），绘制杀伤动力学曲线，从而更精确地区分不同处理的效应窗口。
  \item 统一混样、贴壁与肽脉冲的操作节奏，尽量在同一批次内完成对照与敲除组检测，并在分析中显式纳入批次作为随机效应或分层因素。
\end{enumerate}

\subsection{结论}
综上，本实验在完成 CRISPR/Cas9 介导的 \textit{B2m} 靶向操作与体外共培养杀伤评估流程的基础上，观察到 \textit{B2m} 敲除组在 T 细胞压力下相对存活增加的方向性变化，这与 \textit{B2m} 作为 MHC-I 复合体关键组分、影响抗原呈递与 CTL 识别的机制一致。然而，受限于敲除效率验证不足、孔间方差较大与重复数有限，当前数据尚不足以支持强结论。后续通过补充敲除效率定量、增加重复与优化读出指标，可进一步提高对 \textit{B2m} 在 T 细胞介导杀伤中作用的评估精度。